\documentclass[12pt,a4paper]{article}
\usepackage[utf8]{inputenc}
\usepackage[T1]{fontenc}
\usepackage{geometry}
\geometry{top=2.5cm, bottom=2.5cm, left=2.5cm, right=2.5cm}
\usepackage{times}
\usepackage{hyperref}
\hypersetup{colorlinks=true, linkcolor=blue, citecolor=blue, urlcolor=blue}
\usepackage{xurl}

\setlength{\parindent}{0pt}
\setlength{\parskip}{0.5em}

\begin{document}

% -------------------------------------------------------
% TITLE
% -------------------------------------------------------
\begin{center}
{\Large \textbf{Explainable Hotel Booking Cancellation Prediction Models: Towards Trustworthy Survival Analysis in Hospitality}}
\end{center}

\vspace{1em}

% -------------------------------------------------------
% AUTHOR INFORMATION
% -------------------------------------------------------
\begin{center}
\textbf{Mustafa Cavus}$^{1,*}$
\end{center}

\begin{center}
$^{1}$ Department of Statistics, Eskisehir Technical University, 26555 Eskisehir, T\"{u}rkiye
\end{center}

\begin{center}
$^{*}$ \textbf{Corresponding author:} \href{mailto:mustafacavus@eskisehir.edu.tr}{mustafacavus@eskisehir.edu.tr}
\end{center}

\begin{center}
\textbf{ORCID:} \href{https://orcid.org/0000-0002-6172-5449}{0000-0002-6172-5449}
\end{center}

\vspace{1em}
\hrule
\vspace{1em}

% -------------------------------------------------------
% ABSTRACT
% -------------------------------------------------------
\noindent \textbf{Abstract}

\noindent Hotel booking cancellations impose significant operational and financial risks on revenue management systems, traditionally addressed through static binary classification models. Such approaches fail to capture the temporal dynamics of cancellation behaviors---specifically, \textit{when} a cancellation is most likely to occur between the booking date and check-in. This study addresses this critical gap by formulating cancellation prediction as a survival analysis problem. We benchmark a diverse ensemble of machine learning survival models---including Cox Proportional Hazards, Random Survival Forests, Oblique Random Survival Forests, and Conditional Survival Forests---across two distinct hotel booking datasets with contrasting structural feature spaces. Our findings demonstrate that while standard Random Survival Forests achieve high discriminative rankings, Oblique Random Survival Forests and Conditional Survival Forests deliver superior probability calibration. Furthermore, to advance trustworthy artificial intelligence in hospitality, we integrate time-dependent XAI techniques---Survival Permutation Feature Importance and Partial Dependence Survival Profiles. These tools uncover the non-linear, time-varying influence of core attributes such as \textit{lead time}, \textit{average daily rate}, and \textit{deposit policy} on cancellation risk curves. Specifically, the explanations expose an accelerating cancellation hazard close to check-in among premium-priced bookings driven by guest price sensitivity and intertemporal arbitrage, while isolating a structural anomaly where non-refundable reservations display paradoxically depressed survival rates due to speculative group allocations. The resulting framework equips revenue managers with transparent decision intelligence to dynamically optimize overbooking limits, calibrate cancellation penalties, and mitigate revenue leakage under demand uncertainty.

\vspace{1em}

% -------------------------------------------------------
% KEYWORDS
% -------------------------------------------------------
\noindent \textbf{Keywords:} hotel booking cancellation, XAI, survival analysis, random survival forests, revenue management, trustworthy AI

\vspace{1em}
\hrule
\vspace{1em}

% -------------------------------------------------------
% STATEMENTS AND DECLARATIONS
% -------------------------------------------------------
\noindent \textbf{Statements and Declarations}

\vspace{0.5em}

\noindent \textbf{Competing interests}

\noindent The author has no competing interests to declare that are relevant to the content of this article.

\end{document}